\begin{abstract}
This document proposes four experimental research programs in heterogeneous catalysis, ranked for development, and briefly records six reserve ideas. It reports no new experiments. Its results are a check of each idea against prior work, calculations from published data that size what is at stake, and experimental designs whose rules for expanding or stopping are fixed in advance.

\emph{(1) Physical water management.} Promotion by physically mixed hydrophobic polymers is established. On a Co/SiO\textsubscript{2} Fischer--Tropsch catalyst, a small polydivinylbenzene (PDVB) addition kept CO conversion near 44\% for 1200\,h while a quartz-diluted reference declined. From the published data we calculate a C\textsubscript{5+} output proxy over 25--585\,h of 1.89 times the reference (1.80 per unit catalyst-plus-polymer mass). We also show with a simple storage model that faster outlet-water clearance alone cannot show lower water exposure at active sites. New here: does PDVB added to a catalyst that has already run without it protect the catalyst against a later water challenge, and, conditionally, do measured responses predict output during a small reversible humidity pulse?

\emph{(2) Cyclic oxides.} Li\textsubscript{2}CO\textsubscript{3}-coated La\textsubscript{0.8}Sr\textsubscript{0.2}FeO\textsubscript{3} is a selective oxygen carrier for chemical-looping oxidative dehydrogenation of ethane and was cycled for 1000\,h with Ar purges. The accompanying process model replaces those purges with steam, which was not tested on the coating; \COtwo{} regeneration of carbonate melts has close precedents. New here: does steam in the post-oxidation purge change ethylene output, and is \COtwo{} supplied with the steam (protection) better than \COtwo{} supplied afterwards (recovery)? As an illustration, fully recovering the published Ar-aged C\textsubscript{2+} yield loss would pay for at most 2.02 extra minutes per 16-minute cycle.

\emph{(3) Polymer ethenolysis.} One fresh Na/alumina addition restores the output of a reused Na/alumina--W/silica mixture that converts polyethylene to propylene; pressure effects and chain-population models are also prior work. New here: can a lower ethylene partial pressure let a smaller fresh Na/alumina addition give the same output, and can that saving be predicted? At equal output, halving it would use 25\% less fresh Na/alumina and 16.7\% less total fresh catalyst solid.

\emph{(4) Promoted Ag.} Better selectivity retention by Ni-containing Cs/Re-promoted Ag was disclosed in 1995, confounded by Cs and treatment differences, and chloride adjustment alone changes selectivity and output. New here: a composition-controlled test of when Ni adds ethylene oxide output that chloride operation cannot replace, and a prediction of when a chloride recovery step pays. At fixed output, raising selectivity from 0.90 to 0.91 would cut reaction-level oxygen use by 4.40\% and \COtwo{} formation by 10.99\%.

The rank reflects scientific value: water management and cyclic oxides are the first and second priorities, although polymer ethenolysis and promoted Ag may be quicker first experiments where equipment exists. A program expands only if its measurements predict useful output under a condition withheld from fitting.
\end{abstract}
\clearpage
\tableofcontents
\clearpage

\section*{Overview and recommended priorities}
\addcontentsline{toc}{section}{Overview and recommended priorities}

\paragraph{Purpose and scope.}
This document recommends four research programs for development, in the order of Table~\ref{tab:overview-ranking}. Each links an important catalytic limitation to a specific intervention, a fair material or operating comparator, and a prediction that independent experiments can test. Four is an editorial choice: it gives full treatment to the two strongest development candidates and to two narrower studies that each end in a concrete decision about resource use. Six reserve ideas, whose current limitations do not justify the same experimental scope, are summarized in Section~\ref{sec:shortlist}.

\begin{table}[htbp]
\centering
\footnotesize
\renewcommand{\arraystretch}{1.2}
\caption{The four programs in rank order: what prior work already shows, the open question each addresses, and the reason for its rank.}
\label{tab:overview-ranking}
\begin{tabularx}{\textwidth}{@{}>{\raggedright\arraybackslash}p{0.13\textwidth}>{\raggedright\arraybackslash}X>{\raggedright\arraybackslash}X>{\raggedright\arraybackslash}p{0.2\textwidth}@{}}
\toprule
Program & Already known & Open question (contribution sought) & Reason for rank\\
\midrule
1. Physical water management (Section~\ref{sec:water}) & Physically mixed hydrophobic polymer promotes CO and \COtwo{} hydrogenation.\cite{fang2022,li2023} With a small PDVB addition, Co/SiO\textsubscript{2} kept about 44\% CO conversion in Fischer--Tropsch synthesis for 1200\,h; a quartz-diluted reference started higher but declined strongly.\cite{fang2026} & Does PDVB added after common conditioning without it protect against a later water challenge at one catalyst state? Conditionally, do measured responses predict useful output during a small reversible humidity pulse? & A large published formulation effect. The pulse prediction would not establish a damage law across histories.\\
2. Cyclic oxides (Section~\ref{sec:cyclic}) & Li\textsubscript{2}CO\textsubscript{3}-coated La\textsubscript{0.8}Sr\textsubscript{0.2}FeO\textsubscript{3} (LSF) gives selective chemical-looping oxidative dehydrogenation of ethane and was cycled for 1000\,h with Ar purges.\cite{gao2020,brody2022} Carbonate/hydroxide response to steam and \COtwo{} regeneration have close precedents.\cite{barckholtz2021,chacko2025thesis,fereres2018} & Does steam in the post-oxidation purge change ethylene output, and is \COtwo{} supplied with the steam (protection) better than \COtwo{} supplied afterwards (recovery)? & A specific gap between an Ar-purged durability test and a steam-based process proposal.\\
3. Polymer ethenolysis (Section~\ref{sec:polymer}) & Na/alumina with W/silica converts polyethylene to propylene; one fresh Na/alumina addition restores output during reuse.\cite[Fig.~4B]{conk2024} Pressure effects and chain-population models are prior work.\cite{wang2022thesis,guironnet2020,chen2024} & Can one lower ethylene partial pressure replace part of the fresh Na/alumina while keeping polymer-derived output, and can the saving be predicted? & A concrete catalyst-use decision, although rescue and pressure effects have close prior art.\\
4. Promoted Ag (Section~\ref{sec:silver}) & Ni-containing promoted Ag retained selectivity better than a Ni-free parent in a 1995 patent, with Cs and treatment differences between them.\citep{kemp1995} Dilute Ni suppresses combustion more than ethylene oxide formation on Ag lacking the full Cs/Re promoter system, and chloride operation strongly changes selectivity and rate.\citep{jalil2025,iyer2021,santos2024} & When does Ni add useful ethylene oxide output that chloride adjustment cannot replace, and can that be predicted? & Historical Ni-retention evidence leaves a demanding, narrow material and operating study.\\
\bottomrule
\end{tabularx}
\end{table}

\paragraph{Main results of this analysis.}
The following calculations and corrections are this document's own results; they use published data and are not new measurements.
\begin{itemize}
 \item \emph{Water management.} Integrating Fang et al.'s published conversion and selectivity series over 25--585\,h gives a C\textsubscript{5+} output proxy 1.89 times the reference, or 1.80 per unit catalyst-plus-polymer mass (Eq.~\ref{eq:water-output}). This is not a carbon balance or a lifetime extrapolation. A two-pool storage model shows that outlet water transients cannot locate where water is produced or held, so faster clearance alone does not show drier active sites.
 \item \emph{Cyclic oxides.} As an illustration, fully restoring the published Ar-aged C\textsubscript{2+} yield (47.24\% to 53.2\%) would pay for at most 2.02 additional minutes per 16-minute cycle, before gas and heat costs (Eq.~\ref{eq:cyclic-threshold}). This bounds the value of any recovery step; it does not predict steam damage.
 \item \emph{Polymer ethenolysis.} At equal output, a one-time 0.2\,g makeup would use 25\% less fresh Na/alumina and 16.7\% less total fresh catalyst solid than the published 0.4\,g addition (Table~\ref{tab:polymer-makeup}).
 \item \emph{Promoted Ag.} If ethylene oxide and complete combustion are the only carbon outlets, raising selectivity from 0.90 to 0.91 at fixed ethylene oxide output reduces reaction-level ethylene use by 1.10\%, oxygen use by 4.40\% and \COtwo{} formation by 10.99\%.
 \item \emph{Withdrawn claims.} A faster outlet-water transient does not by itself show lower active-site water exposure; fresh-Na rescue does not identify selective loss of initiation; Ni compatibility and retention on promoted Ag are prior art; and published mannose results contradict treating mannose as an inactive reservoir in tungsten sugar conversion (Section~\ref{sec:shortlist}).
\end{itemize}

\paragraph{Priorities and feasibility.}
The ranking concerns scientific development and space in the document, not a plan to run four campaigns at once. Water management is the first priority and cyclic oxides the second. Polymer ethenolysis and promoted Ag may nevertheless offer the best first experiment where their reactors, reference materials and analytical methods already exist. Laboratory access, preparation reproducibility, achievable uncertainty and actual costs have not been established. None of the four yet justifies an open-ended campaign.

\paragraph{Three kinds of statement.}
The document keeps three kinds of statement separate. \emph{Published observations} are attributed to their sources and keep the sources' preparation, measurement and time limits. \emph{Calculations} are deductions from stated assumptions or identified published data; they are not new measurements. \emph{Proposed experiments} have not been performed; each states the information a decision needs, the possible outcomes and the limits of interpretation. The contribution is a testable research design, not advance confirmation of a mechanism. In particular, the document does not establish lower local water exposure at Co sites, steam-induced Li loss, selective loss of polymer initiation, or preservation of a particular Ag--Re oxygen-transfer state.

\paragraph{When to expand or stop a program.}
Each design compares against a fair comparator and a simpler credible operating alternative: PDVB versus quartz, both added after common conditioning; ordinary steam versus Ar before any timed \COtwo{} treatment; ethylene pressure crossed with fresh-Na dose; and Ni-treated versus sham-processed Ag under ordinary and recovery policies, plus untreated Ag under ordinary operation. A program expands only when its measurements identify a material or operating choice and correctly predict that choice's response under a condition or history withheld from fitting. Further characterization is justified only when it separates explanations that would lead to different decisions. Output accounting stays complete even when a mechanistic interpretation succeeds.

\paragraph{Design rules common to all programs.}
The chapters apply the following rules without restating them each time.
\begin{itemize}
 \item \emph{Replication.} Replicates are independent catalyst charges or specimens, randomized within preparation and run-order blocks. Repeated analyses or cycles on one bed are not replicates. Replicate and cycle counts come from pilot variability and the smallest important difference, both fixed before the decisive runs.
 \item \emph{Null results.} A null result counts as evidence against an important effect only if its uncertainty interval (including preparation, run, calibration and balance uncertainty) excludes the smallest important difference. Positive and null results alike hold only for the tested material, conditions, duration and history; no finite laboratory campaign here can establish lifetime tolerance or competitive industrial lifetime.
 \item \emph{Predictions.} Models, predictions and their uncertainty intervals are fixed before the withheld test, which uses independent charges. A failed prediction is reported as a failure, not refitted. Where a chapter names a simpler baseline, a more detailed model or added mechanism is credited only if it predicts the withheld result better than that baseline; otherwise only the simpler rule is claimed.
 \item \emph{Accounting.} Practical output comparisons use complete elapsed time over a horizon fixed in advance, counting every treatment, handling and recovery step that differs between arms, with all catalyst, gas and product inventories closed. Absolute outputs are reported alongside ratios or fractions.
 \item \emph{Ex situ states.} Characterization after quench or transfer is interpreted only after showing that sampling does not change the measured property; otherwise the attribution is withheld.
\end{itemize}
